\section{Modelo de IA e implantação no firmware}
\label{sec:modelo}

\subsection{Treinamento no computador e inferência no dispositivo}

Random Forest combina diferentes árvores treinadas com aleatorização
de amostras e atributos \cite{breiman2001}. No projeto, o treinamento
foi implementado com scikit-learn \cite{pedregosa2011}. A configuração
compacta tem 32 árvores, profundidade máxima 6, mínimo de cinco amostras
por folha, balanceamento \texttt{balanced\_subsample} e semente 42.
Esses parâmetros e o relatório por classe constam do artefato técnico
\path{docs/embedded-ai-metrics.json}.

A Figura~\ref{fig:treinamento} diferencia treinamento, avaliação e
exportação. O script \path{ai/train_embedded.py} gera o cabeçalho
\path{firmware/ms_monitor/risk_model.h}. Os parâmetros aprendidos são
representados por tabelas constantes de nós e raízes; não são limiares
manuais definidos para inferir presença de fungos.
O arquivo \texttt{.joblib} é mantido no computador para reprodução e
comparação e não é carregado pela ESP32 nem pela API atual.

\begin{figure}[htbp]
\centering
\caption{Preparação do modelo embarcado e verificação da exportação}
\label{fig:treinamento}
\resizebox{0.96\textwidth}{!}{\begin{tikzpicture}[x=1cm,y=1cm]
\node[block,text width=8cm] (data) at (0,0)
{Base local: 43.200 registros / 20 lotes\\Três medições e diferença térmica};
\node[block,text width=8cm] (split) at (0,-1.6)
{Separação por lote\\Primeira divisão de cinco partições};
\node[block,text width=5.4cm] (train) at (-3.6,-3.5)
{Treinamento: 15 lotes\\Random Forest compacto\\32 árvores / profundidade até 6};
\node[block,text width=5.4cm] (test) at (3.6,-3.5)
{Teste separado: 5 lotes\\10.800 registros\\Não reutilizados no ajuste};
\node[block,text width=5.4cm] (export) at (-3.6,-5.7)
{Exportar parâmetros aprendidos\\Cabeçalho C++ e versão do modelo};
\node[block,text width=5.4cm] (metrics) at (3.6,-5.7)
{Avaliar previsões\\Métricas por classe e\\limites de representatividade};
\node[block,text width=5.4cm] (reference) at (3.6,-7.7)
{Preditor Python de referência\\Aplicar ao dataset completo};
\node[block,text width=11.5cm] (parity) at (0,-9.7)
{Comparar o preditor C++ com o Python no dataset completo\\
Paridade de implementação não equivale a validação agronômica};
\draw[flow] (data) -- (split);
\draw[flow] (split.south) -- ++(0,-0.35) -| (train.north);
\draw[flow] (split.south) -- ++(0,-0.35) -| (test.north);
\draw[flow] (train) -- (export);
\draw[flow] (test) -- (metrics);
\draw[flow] (train.east) -- (test.west) node[midway,above,edge label] {modelo};
\fill (0,-3.5) circle (1.3pt);
\draw[flow] (0,-3.5) -- (0,-7.7) -- (reference.west);
\draw[flow] (export.south) -- (-3.6,-8.7) -| ([xshift=-2cm]parity.north);
\draw[flow] (reference.south) -- ++(0,-0.35) -| ([xshift=2cm]parity.north);
\end{tikzpicture}
}
\source{elaboração própria a partir dos scripts de treinamento e verificação}
\end{figure}

Os dados locais são descritos na documentação como sintéticos: 43\,200
registros distribuídos em 20 lotes. Não há, no material examinado,
comprovação de que os rótulos derivem de análises microbiológicas
dos grãos do protótipo. A identificação do conjunto é preservada
por seu hash SHA-256 no relatório de treinamento.

O código usa a \textbf{primeira divisão} produzida por
\texttt{StratifiedGroupKFold} com cinco partições, mantendo lotes
inteiros separados. O experimento utiliza 32\,400 registros de 15
lotes no treinamento e 10\,800 registros de cinco lotes no teste.
Não se trata de uma média de desempenho de cinco rodadas de validação
cruzada, nem de uma validação temporal independente.

\subsection{Saídas do classificador e seu significado}

A execução retorna um dos rótulos \texttt{low}, \texttt{moderate},
\texttt{high} ou \texttt{critical}. Neste relatório, eles são
interpretados como níveis nominais do conjunto de treinamento,
respectivamente baixo, moderado, alto e crítico.
Esses nomes não estabelecem critérios agronômicos validados nem
recomendações automáticas de intervenção.

Para cada classe $c$, o preditor calcula a média das probabilidades
das folhas atingidas pelas $M=32$ árvores:
\begin{equation}
p_c(\mathbf{x})=\frac{1}{M}\sum_{m=1}^{M}p_{m,c}(\mathbf{x}),
\qquad
\widehat{c}=\arg\max_c p_c(\mathbf{x}).
\end{equation}
A confiança reportada é $\max_c p_c(\mathbf{x})$; não corresponde a
uma medida experimental de certeza sobre a condição sanitária do lote.
O índice apresentado em \texttt{risk\_score} segue a convenção
\begin{equation}
s=0\,p_{\mathrm{low}}+0{,}33\,p_{\mathrm{moderate}}
+0{,}66\,p_{\mathrm{high}}+1\,p_{\mathrm{critical}}.
\label{eq:score}
\end{equation}
Os pesos são uma escala de apresentação adotada no código. Mesmo
quando expresso como percentual no dashboard, $s$ não é uma
probabilidade calibrada de haver fungos ou micotoxinas.

\subsection{Execução e independência da rede}

O fluxo da Figura~\ref{fig:algoritmo} segue a ordem do firmware.
O tempo de inferência é medido ao redor da chamada local do preditor
com \texttt{micros()}, antes de verificar a conectividade.
Sem Wi-Fi, o dispositivo continua lendo os sensores, classificando
entradas válidas e registrando o resultado na serial. Os dados
desse período não são enfileirados para envio posterior.

Quando há rede, ocorre uma tentativa de POST por ciclo. O firmware
não testa antecipadamente a disponibilidade da API: interpreta a
resposta HTTP após tentar o envio. O código considera 201 como
aceitação e registra os demais códigos ou erros. Os limites
configurados são 700~ms para conexão e 1000~ms para leitura HTTP;
não representam garantia de latência total do ciclo.

\begin{figure}[p]
\centering
\caption{Fluxograma da aquisição e da inferência no firmware}
\label{fig:algoritmo}
\resizebox{0.94\textwidth}{!}{\begin{tikzpicture}[x=1cm,y=1cm]
\node[terminal,text width=3cm] (start) at (0,0) {Início};
\node[block,text width=9cm] (init) at (0,-1.2)
{Inicializar sensores, interrupção e Wi-Fi\\Modelo já incluído no firmware};
\node[block,text width=8cm] (wait) at (0,-2.8)
{Aguardar ciclo nominal de 2,5 s};
\node[block,text width=10cm] (read) at (0,-4.4)
{Fechar janela do contador em seção crítica\\Ler DO, DHT22 e MAX6675; validar cada sensor};
\node[decision,text width=4cm] (valid) at (0,-6.6)
{DHT22 e sonda\\válidos?};
\node[block,text width=5.2cm] (invalid) at (-3.7,-8.8)
{Não executar o modelo\\Classe unavailable\\Score e confiança null};
\node[block,text width=5.2cm] (infer) at (3.7,-8.8)
{Preparar quatro atributos\\Executar Random Forest\\Medir tempo de inferência};
\node[block,text width=9.5cm] (serial) at (0,-10.7)
{Registrar medições e classificação na serial\\Preservar telemetria válida de vibração};
\node[decision,text width=4cm] (wifi) at (0,-12.9)
{Wi-Fi\\conectado?};
\node[block,text width=5.2cm] (offline) at (-3.7,-15.1)
{Sem envio e sem fila\\Reconexão temporizada\\Leituras locais continuam};
\node[block,text width=5.2cm] (post) at (3.7,-15.1)
{Montar JSON e tentar POST\\Timeout finito\\Registrar HTTP 201 ou falha};
\node[terminal,text width=5cm] (next) at (0,-17.1) {Próximo ciclo};
\draw[flow] (start) -- (init);
\draw[flow] (init) -- (wait);
\draw[flow] (wait) -- (read);
\draw[flow] (read) -- (valid);
\draw[flow] (valid.west) -| (invalid.north) node[pos=0.25,above,edge label] {não};
\draw[flow] (valid.east) -| (infer.north) node[pos=0.25,above,edge label] {sim};
\draw[flow] (invalid.south) -- ++(0,-0.3) -| ([xshift=-2cm]serial.north);
\draw[flow] (infer.south) -- ++(0,-0.3) -| ([xshift=2cm]serial.north);
\draw[flow] (serial) -- (wifi);
\draw[flow] (wifi.west) -| (offline.north) node[pos=0.25,above,edge label] {não};
\draw[flow] (wifi.east) -| (post.north) node[pos=0.25,above,edge label] {sim};
\draw[flow] (offline.south) -- ++(0,-0.3) -| ([xshift=-1.5cm]next.north);
\draw[flow] (post.south) -- ++(0,-0.3) -| ([xshift=1.5cm]next.north);
\draw[flow] (next.west) -- (-7.3,-17.1) -- (-7.3,-2.8) -- (wait.west);
\end{tikzpicture}
}
\source{elaboração própria a partir de \texttt{ms\_monitor.ino}; a rotina de interrupção do SW-420 opera em paralelo ao ciclo principal}
\end{figure}

O modelo embarcado é identificado como \modelversion. O relatório de
compilação registrado no projeto, para core Espressif 3.3.12 e
FQBN \path{esp32:esp32:heltec_wifi_lora_32_V2}, informa
1\,043\,428 bytes de programa, diante de um limite de 3\,342\,336 bytes,
e 48\,824 bytes de variáveis globais na região de 327\,680 bytes
contabilizada pela ferramenta. As tabelas de nós e raízes ocupam
102\,496 bytes no compilador nativo de verificação.

Esses números sustentam a análise de compatibilidade da compilação,
mas não medem o heap livre em execução, a duração da inferência
na placa ou o consumo elétrico. Tampouco se deve confundir a região
de memória descrita pelo compilador com a SRAM física total do chip.
\FloatBarrier
